%========================================================== GATE IDENTITY
\section{Knowing which gate is which}\label{sec:identity}

\subsection{The problem, and why it is not a perception problem}

The policy is told which gate it is flying at. In the simulator that arrives for
free. On the real course nothing publishes it, ten gates are identical
\SI{2.7}{m} squares, and \textbf{we asked the organizers whether we may mark them
and the answer was no}.

Measured on our own course: \textbf{34\,\% of frames have two or more gates in
view at once.} So ``look at the gate and recognise it'' is not available from
appearance, ever.

\begin{explain}[The reframe that makes this tractable]
Gate identity is not a perception problem. It is a \textbf{localisation}
problem. If you know where you are, identity is a lookup: project all ten gates
from the map and see which one you are pointed at. If you do not know where you
are, no amount of staring at an identical square will tell you.

Every published system reaches the same conclusion and none of them tries to
recognise gates individually.
\end{explain}

\subsection{What the policy actually reads}\label{sec:statevector}

Worth stating explicitly, because the argument below turns on it. Each frame is
51 numbers, and the policy reads \textbf{32 frames of history}, so 1632 inputs.

\begin{table}[H]\centering\footnotesize
\begin{tabularx}{\linewidth}{@{}p{1.6cm} p{3.1cm} X@{}}
\toprule
Index & Name & Meaning \\
\midrule
0--15 & \code{u0,v0 \dots u7,v7} & \textbf{Eight gate corners} as image positions, normalised by 640$\times$360. \textbf{$-1.0$ in both when not seen} \\
16--23 & \code{vis0 \dots vis7} & Is that corner seen? 1.0 or 0.0 \\
24--25 & \code{roll}, \code{pitch} & Attitude from the flight controller, radians \\
26--28 & \code{gx, gy, gz} & Gyro rates, body NED, clipped $\pm$8\,rad/s \\
29--31 & \code{vx, vy, vz} & Commanded body velocity, NED, clipped $\pm$20\,m/s \\
32--49 & \code{gate0 \dots gate17} & \textbf{Which gate is next} --- one slot set to 1.0 \\
50 & \code{lap\_frac} & Progress round the lap, $\mathrm{index}/17$ \\
\bottomrule
\end{tabularx}
\caption{The 51-number observation frame. Verified against the training code element by element (Section~\ref{sec:built}).}
\end{table}

The eight corners are \textbf{one gate only} --- the one being flown at, not all
ten. Corners 0--3 are the outer \SI{2.7}{m} frame, clockwise from top-left;
corners 4--7 are the inner \SI{1.5}{m} opening in the same order. Having both
rings is what encodes \emph{range}: their true sizes are known, so apparent size
gives distance.

Three things the policy is \textbf{not} given, which matter for what follows:
no distance, no position or altitude or heading (roll and pitch but not yaw),
and no sight of the \emph{next} gate. Everything beyond the gate in view comes
from the one-hot label and what the network memorised in training.

\begin{warnbox}[The consequence]
Nothing in indices 0--31 distinguishes gate~3 from gate~8. All ten are identical
squares. The \textbf{only} channel carrying course position is the one-hot in
32--49 --- which is 576 of the 1632 numbers once the 32 history frames are
counted, over a third of everything the network reads.
\end{warnbox}

\subsection{``The model will learn the order from simulation''}

A reasonable-sounding proposal was put forward: label the gates in simulation,
train on that, and the model will know the order on the real unlabelled gates.

The mechanical objection is that the label is an \textbf{input}, not an output.
Training with it does not teach the network the order; it teaches the network to
\emph{depend on being told} the order. At run time something must still supply
those 18 numbers, and a network cannot supply its own input.

But that is an argument, and arguments lose to measurements. So we measured it.

\paragraph{The experiment.} Gene's model, flown in \emph{his} simulator on
\emph{his} track --- his gate positions, his \SI{90}{\degree} camera, the
\SI{607.6}{g} drone --- the conditions where it works well. 512 drones at once,
\SI{20}{s} each. Identical physics, course, camera and seed across all five runs.
\textbf{The only thing changed is where the gate label comes from.}

\begin{table}[H]\centering\footnotesize
\begin{tabular}{@{}lrrrr@{}}
\toprule
Label source & \makecell{Gates per\\crash} & \makecell{Crashes /\\100 gates} & \makecell{Median\\lap} & \makecell{Laps\\completed} \\
\midrule
\textbf{Correct, as trained} & \textbf{26.9} & \textbf{3.4} & \textbf{7.45\,s} & \textbf{515} \\
Stuck on the start gate & 1.8 & 51.5 & --- & 0 \\
No label at all & 1.1 & 85.6 & --- & 0 \\
Random label each frame & 1.1 & 72.9 & --- & 0 \\
Off by one & 1.0 & 80.7 & --- & 0 \\
\bottomrule
\end{tabular}
\caption{The gate label is load-bearing. Raw data in \nolinkurl{~/aigp/eval/ctx/}; reproduce with \nolinkurl{run_ctx.sh}.}\label{tab:ctx}
\end{table}

\begin{verdict}[What this settles]
\begin{itemize}
\item \textbf{With the correct label the model is healthy}: \SI{7.45}{s} laps,
\SI{16}{m/s} at the gate, 515 complete laps, about 27 gates between crashes.
Exactly the performance it always had.
\item \textbf{Without a label it completes zero laps.} Not one, across any of the
four broken conditions. It crashes after roughly one gate.
\item ``Train with labels in simulation, then have nothing supplying them in
reality'' \emph{is} the ``no label at all'' row.
\item \textbf{Off-by-one is as bad as no label at all} (1.0 vs 1.1). This is not
only an argument for having a gate counter --- it is an argument for the counter
being \emph{right}. A counter that slips once is worth nothing.
\item \textbf{A frozen counter is the least-bad failure} (1.8). So when
confidence is lost, hold the last index rather than guess a new one. That is how
ours is built.
\end{itemize}
\end{verdict}

The instinct behind the proposal is not wrong, and is a real published
technique --- but it requires \emph{changing how we train}, not merely training
with labels present. Either drop the input entirely and force the network to
infer position from its 32 frames of history, or train a teacher with the
privileged label and distil it into a student that sees only what the drone
sees. Both work. Both have the same weakness: when the network's internal guess
is wrong there is no error signal, so it is confidently wrong. The approach
below at least knows when it is lost.

\subsection{What we build instead}

The published systems --- Swift\X, AlphaPilot\X\ and the 2026 dual pose-graph
work\X\ --- all share one shape: detect gate corners, solve a relative pose per
gate, match against the \emph{known map} in metric space, and fuse into a filter.
Swift states it plainly: each observation is assigned to ``the nearest gate in
the known trajectory layout''.

Two things we do differently, both because we measured rather than copied.

\paragraph{No visual-inertial odometry.} Every one of those systems leans on VIO.
We have one forward camera and a \SI{25}{W} budget. So we measured whether we
need it: a gate is in view \textbf{75--79\,\%} of frames, the worst gap is
\SI{1.5}{s} at \SI{5}{m/s}, and dead reckoning drifts \textbf{under a metre}
across that gap --- against gates that are \SI{7.54}{m} apart. VIO would be
solving a problem we do not have.

\paragraph{Throw away the PnP rotation.} Planar PnP is ill-conditioned near
head-on, which is where a racing drone spends its life. But we do not need its
rotation: the flight controller already gives us attitude, far more accurately.

\begin{table}[H]\centering\footnotesize
\begin{tabular}{@{}lrrr@{}}
\toprule
Corner noise & Range & Full 6-DOF PnP & \textbf{PnP translation + FC attitude} \\
\midrule
2\,px & 8\,m & 2.12\,m & \textbf{0.09\,m} \\
5\,px & 8\,m & 3.15\,m & \textbf{0.29\,m} \\
10\,px & 12\,m & 6.81\,m & \textbf{1.38\,m} \\
\bottomrule
\end{tabular}
\caption{Discarding the ill-conditioned half of the PnP solution is worth 23$\times$ at the operating point. This is not from the literature; it fell out of the accuracy test.}\label{tab:pnp}
\end{table}

\subsection{The pipeline}

\begin{enumerate}
\item \textbf{Detector} --- finds a gate's eight corners in the image.
\emph{Still to be built; needs real footage.}
\item \textbf{PnP} --- corners plus the known \SI{2.7}{m}/\SI{1.5}{m} geometry
give range and bearing to the gate. Planar PnP returns two mirrored solutions;
we pick by reprojection error and break ties with \textbf{gravity}, since gates
stand vertically and we know which way is down.
\item \textbf{Position} --- range and bearing, plus the flight controller's
attitude, plus the gate's known map position, give \emph{where we are}.
\item \textbf{Filter} --- carries position through the gaps by integrating thrust
and attitude, corrected by each fix. It also estimates \textbf{thrust scale},
which is what makes a barometer unnecessary --- and the barometer is unusable
anyway with the props spinning.
\item \textbf{Tracker} --- knowing where we are, project all ten gates and see
which one is ahead. Advance on a confirmed crossing. Fill the 18 one-hot slots.
\end{enumerate}

In one line: \emph{the gate looks this big and sits here in frame} $\rightarrow$
\emph{so it is \SI{8.3}{m} away, \SI{12}{\degree} left} $\rightarrow$ \emph{the
map says that gate is at (12.19, 45.11)} $\rightarrow$ \emph{so I am at
(10.4, 37.2)} $\rightarrow$ \emph{therefore this is gate 3} $\rightarrow$
\emph{set slot 35 to 1.0}.

\subsection{Evidence}\label{sec:built}

\begin{table}[H]\centering\footnotesize
\begin{tabularx}{\linewidth}{@{}p{4.3cm} X@{}}
\toprule
Check & Result \\
\midrule
Observation frame vs training code & \textbf{Exactly zero} difference, 51 numbers $\times$ 200 cases \\
History stack vs training code & \textbf{Exactly zero} difference across 1632 numbers \\
Projection vs training code & $3.9\times10^{-6}$\,px, zero visibility mismatches over 3200 corners \\
Policy in NumPy vs PyTorch & $2.5\times10^{-6}$ max error on the real checkpoint, \SI{126}{\micro s} per inference \\
Position, normal conditions & \textbf{0.88\,m median, 2.79\,m p95} against \SI{7.54}{m} gate separation \\
\dots with 20\,\% thrust error & 0.88\,m / 2.79\,m (thrust scale learned 0.895 vs true 0.900) \\
\dots drag model 2$\times$ wrong & 0.80\,m / 3.01\,m \\
\dots 10\,px corner noise & 1.01\,m / 2.38\,m \\
\dots \SI{7}{m} detector range & 0.77\,m / 2.29\,m \\
\dots \SI{12}{m/s} flight & 0.77\,m / 1.57\,m \\
Dead reckoning with no fixes & 88.9\,m / 177\,m --- useless alone, which is the point \\
Gate counting, 17 conditions & 22/22 crossings, zero errors (50\,\% dropout, 20 false positives/s, \SI{100}{ms} latency, 1\,m pose noise) \\
\bottomrule
\end{tabularx}
\caption{Everything above is reproducible with \code{python run\_all.py} in \nolinkurl{flight/}.}
\end{table}

\begin{warnbox}[The one known limit: heading drift]
The position fix is computed \emph{using} our heading, so a wrong heading
corrupts the fix directly and no filter can recover --- the measurement itself is
wrong. Measured tolerance:

\begin{center}\footnotesize
\begin{tabular}{@{}lrrl@{}}
\toprule
Yaw drift & Median & p95 & \\
\midrule
\SI{0.05}{\degree\per\second} & 0.28\,m & 0.97\,m & fine \\
\SI{0.5}{\degree\per\second} & 0.90\,m & 2.81\,m & fine \\
\SI{1.0}{\degree\per\second} & 1.47\,m & 6.81\,m & \textbf{breaks} \\
\bottomrule
\end{tabular}
\end{center}

So it tolerates about \SI{0.5}{\degree\per\second} and fails near
\SI{1}{\degree\per\second}. A calibrated Betaflight gyro should be an order of
magnitude better than that, but this is now a specific number to \textbf{verify
from the measurement flight} rather than assume.
\end{warnbox}
